\documentclass[11pt]{article}
\usepackage[margin=2cm]{geometry}\usepackage[T1]{fontenc}\usepackage{fancyvrb}\usepackage{xcolor}\usepackage{amsmath}
\title{GATE DA 2027 -- Missing Topics Addendum\\\large Ramu and Somu}\date{}
\begin{document}\maketitle
This addendum fills the syllabus topics that the main notes covered only briefly. Every code block was run, and the printed output shown is the real output.
\section{Data Structures and Sorting (Programming, DSA)}
\subsection*{Linked lists}
Ramu: Why not just use an array?\\Somu: A linked list changes at the front in $O(1)$ by moving one pointer, but reaching the $k$-th item costs $O(k)$. Reversal flips each pointer once: $O(n)$.
\begin{Verbatim}[fontsize=\footnotesize,frame=single]
class Node:
    def __init__(self, val, nxt=None):
        self.val, self.nxt = val, nxt      # value + pointer to next node
head = Node(1, Node(2, Node(3)))           # 1 -> 2 -> 3
def to_list(n):
    out = []
    while n: out.append(n.val); n = n.nxt  # walk the pointers: O(n)
    return out
def reverse(head):
    prev = None
    while head:
        nxt = head.nxt                     # remember the rest
        head.nxt = prev                    # flip the pointer
        prev, head = head, nxt             # step forward
    return prev
print(to_list(head), to_list(reverse(head)))
# UPDATE: insert at the head is O(1); finding the k-th item is O(k)
\end{Verbatim}
\textbf{Output:}
\begin{Verbatim}[fontsize=\footnotesize,frame=single]
[1, 2, 3] [3, 2, 1]
\end{Verbatim}
\subsection*{Binary search trees}
Smaller keys go left, bigger go right. Search cost equals the height: $O(\log n)$ balanced, $O(n)$ when keys arrive sorted.
\begin{Verbatim}[fontsize=\footnotesize,frame=single]
class T:
    def __init__(self, k): self.k, self.l, self.r = k, None, None
def insert(t, k):
    if t is None: return T(k)
    if k < t.k: t.l = insert(t.l, k)       # smaller goes left
    else: t.r = insert(t.r, k)             # bigger goes right
    return t
def height(t): return 0 if t is None else 1 + max(height(t.l), height(t.r))
def inorder(t): return inorder(t.l) + [t.k] + inorder(t.r) if t else []
root = None
for k in [50, 30, 70, 20, 40, 60, 80]: root = insert(root, k)
print(inorder(root), height(root))         # in-order = sorted
skew = None
for k in [10, 20, 30, 40, 50, 60, 70]: skew = insert(skew, k)
print(height(skew))                        # sorted input = a chain
# UPDATE: search costs O(height): log n if balanced, n if skewed
\end{Verbatim}
\textbf{Output:}
\begin{Verbatim}[fontsize=\footnotesize,frame=single]
[20, 30, 40, 50, 60, 70, 80] 3
7
\end{Verbatim}
\subsection*{Hash tables and collisions}
A hash function maps a key to a slot. Different keys can share a slot (a collision); chaining keeps a list per slot. Load factor $\alpha=n/m$; the average lookup is $1+\alpha$.
\begin{Verbatim}[fontsize=\footnotesize,frame=single]
M = 5
table = [[] for _ in range(M)]             # each slot holds a chain
def put(k): table[k % M].append(k)         # hash = k mod 5
for k in [12, 22, 7, 17, 3, 27]: put(k)
print(table)
print(len(table[2]))                       # keys 12,22,7,17,27 share slot 2
print(6 / M)                               # load factor = n / slots
# UPDATE: average cost is 1 + load factor; a bad hash makes it O(n)
\end{Verbatim}
\textbf{Output:}
\begin{Verbatim}[fontsize=\footnotesize,frame=single]
[[], [], [12, 22, 7, 17, 27], [3], []]
5
1.2
\end{Verbatim}
\subsection*{Insertion and selection sort}
Insertion sort is $O(n)$ on sorted input and $O(n^2)$ on reversed input. Selection sort always does about $n^2/2$ comparisons but at most $n$ swaps.
\begin{Verbatim}[fontsize=\footnotesize,frame=single]
def isort(a):
    a, shifts = a[:], 0
    for i in range(1, len(a)):
        key, j = a[i], i - 1
        while j >= 0 and a[j] > key:       # shift bigger items right
            a[j + 1] = a[j]; j -= 1; shifts += 1
        a[j + 1] = key
    return a, shifts
def ssort(a):
    a, swaps = a[:], 0
    for i in range(len(a)):
        m = min(range(i, len(a)), key=lambda x: a[x])   # find the smallest
        a[i], a[m] = a[m], a[i]; swaps += 1
    return a, swaps
print(isort([1, 2, 3, 4, 5, 6])[1], isort([6, 5, 4, 3, 2, 1])[1])
print(ssort([5, 2, 4, 1, 3]))
# UPDATE: insertion is O(n) if sorted, O(n^2) if reversed
\end{Verbatim}
\textbf{Output:}
\begin{Verbatim}[fontsize=\footnotesize,frame=single]
0 15
([1, 2, 3, 4, 5], 5)
\end{Verbatim}
\subsection*{Quick sort}
Average $O(n\log n)$; with a first-item pivot on sorted input it degrades to $n(n-1)/2$ comparisons. A random pivot makes the bad case unlikely.
\begin{Verbatim}[fontsize=\footnotesize,frame=single]
import random
def qsort(a, c):
    if len(a) <= 1: return a
    p = a[0]                               # pivot = first item
    lo = [x for x in a[1:] if x < p]
    hi = [x for x in a[1:] if x >= p]
    c[0] += len(a) - 1                     # comparisons this level
    return qsort(lo, c) + [p] + qsort(hi, c)
random.seed(2)
n = 200
c1, c2 = [0], [0]
qsort(random.sample(range(n), n), c1)
qsort(list(range(n)), c2)
print(c1[0], c2[0], n * (n - 1) // 2)      # random vs already sorted
# UPDATE: random pivot keeps the expected cost at n log n
\end{Verbatim}
\textbf{Output:}
\begin{Verbatim}[fontsize=\footnotesize,frame=single]
1707 19900 19900
\end{Verbatim}
\section{Probability and Statistics}
\subsection*{Mean, median, mode, spread}
The mean chases outliers, the median does not. Population standard deviation is the root of the average squared deviation.
\begin{Verbatim}[fontsize=\footnotesize,frame=single]
data = [2, 4, 4, 4, 5, 5, 7, 9]
n = len(data)
mean = sum(data) / n
s = sorted(data)
median = (s[n//2 - 1] + s[n//2]) / 2       # middle of the sorted list
mode = max(set(data), key=data.count)      # most frequent value
var = sum((x - mean)**2 for x in data) / n # population variance
print(mean, median, mode, var ** 0.5)
big = data + [100]                         # add one outlier
bs = sorted(big)
print(round(sum(big)/len(big), 2), bs[len(big)//2])
# UPDATE: the mean chases outliers, the median does not
\end{Verbatim}
\textbf{Output:}
\begin{Verbatim}[fontsize=\footnotesize,frame=single]
5.0 4.5 4 2.0
15.56 5
\end{Verbatim}
\subsection*{Covariance and correlation}
$r=\mathrm{cov}(x,y)/(\sigma_x\sigma_y)$. $r=0$ does not mean independent: $y=x^2$ has covariance $0$.
\begin{Verbatim}[fontsize=\footnotesize,frame=single]
xs = [1, 2, 3, 4, 5, 6]
ys = [2.1, 3.9, 6.2, 7.8, 10.1, 12.2]
n = len(xs)
mx, my = sum(xs)/n, sum(ys)/n
cov = sum((x-mx)*(y-my) for x, y in zip(xs, ys)) / n
sx = (sum((x-mx)**2 for x in xs) / n) ** 0.5
sy = (sum((y-my)**2 for y in ys) / n) ** 0.5
print(round(cov, 3), round(cov / (sx * sy), 4))   # r = cov / (sx sy)
zs = [-3, -2, -1, 0, 1, 2, 3]
zy = [z * z for z in zs]                   # y is fully determined by x
mz, mw = 0, sum(zy) / 7
print(sum((a - mz) * (b - mw) for a, b in zip(zs, zy)) / 7)
# UPDATE: r measures LINEAR link only
\end{Verbatim}
\textbf{Output:}
\begin{Verbatim}[fontsize=\footnotesize,frame=single]
5.892 0.9991
0.0
\end{Verbatim}
\subsection*{The t test}
$t=(\bar x-\mu_0)/(s/\sqrt n)$ with $s$ using $n-1$. Compare $|t|$ with the critical value for $n-1$ degrees of freedom (2.262 for 9 df at 5\% two-sided).
\begin{Verbatim}[fontsize=\footnotesize,frame=single]
data = [5.1, 4.9, 5.6, 5.8, 6.0, 5.5, 5.3, 6.1, 5.7, 5.9]
mu0 = 5.0                                  # claim: the mean is 5.0
n = len(data)
m = sum(data) / n
sd = (sum((x - m)**2 for x in data) / (n - 1)) ** 0.5   # sample sd, n-1
t = (m - mu0) / (sd / n ** 0.5)            # t statistic
crit = 2.262                               # t table, df=9, 5% two-sided
print(round(m, 3), round(sd, 3), round(t, 3))
print(abs(t) > crit)                       # True -> reject the claim
# UPDATE: use t (not z) when sigma is unknown and n is small
\end{Verbatim}
\textbf{Output:}
\begin{Verbatim}[fontsize=\footnotesize,frame=single]
5.59 0.393 4.749
True
\end{Verbatim}
\section{Machine Learning}
\subsection*{Hierarchical clustering and linkage}
Merge the closest clusters repeatedly. Single linkage uses the closest pair (chains), complete linkage the farthest pair (compact).
\begin{Verbatim}[fontsize=\footnotesize,frame=single]
pts = [1, 2, 3, 5, 7, 8, 9]
cl = [[p] for p in pts]                    # start: every point alone
def d_single(a, b): return min(abs(x - y) for x in a for y in b)
def d_complete(a, b): return max(abs(x - y) for x in a for y in b)
def run(dist):
    c, merges = [x[:] for x in cl], []
    while len(c) > 1:
        i, j = min(((i, j) for i in range(len(c)) for j in range(i+1, len(c))),
                   key=lambda ij: dist(c[ij[0]], c[ij[1]]))
        merges.append((c[i] + c[j], dist(c[i], c[j])))   # merge the closest pair
        c = [x for k, x in enumerate(c) if k not in (i, j)] + [c[i] + c[j]]
    return merges
print(run(d_single)[-2:])
print(run(d_complete)[-2:])
# UPDATE: single linkage chains; complete linkage keeps clusters compact
\end{Verbatim}
\textbf{Output:}
\begin{Verbatim}[fontsize=\footnotesize,frame=single]
[([5, 3, 1, 2], 2), ([9, 7, 8, 5, 3, 1, 2], 2)]
[([1, 2, 3, 5], 4), ([9, 7, 8, 1, 2, 3, 5], 8)]
\end{Verbatim}
\subsection*{Linear discriminant analysis}
Fisher direction $w=S_w^{-1}(m_1-m_0)$, supervised, assumes equal class covariance.
\begin{Verbatim}[fontsize=\footnotesize,frame=single]
x0 = [(1, 2), (2, 3), (3, 3), (2, 1)]      # class 0
x1 = [(6, 5), (7, 7), (8, 6), (7, 4)]      # class 1
mean = lambda xs: [sum(c) / len(xs) for c in zip(*xs)]
m0, m1 = mean(x0), mean(x1)
def scatter(xs, m):
    s = [[0, 0], [0, 0]]
    for x in xs:
        d = [x[0] - m[0], x[1] - m[1]]
        for i in range(2):
            for j in range(2): s[i][j] += d[i] * d[j]
    return s
S0, S1 = scatter(x0, m0), scatter(x1, m1)
Sw = [[S0[i][j] + S1[i][j] for j in range(2)] for i in range(2)]  # within-class
det = Sw[0][0] * Sw[1][1] - Sw[0][1] * Sw[1][0]
inv = [[Sw[1][1] / det, -Sw[0][1] / det], [-Sw[1][0] / det, Sw[0][0] / det]]
dm = [m1[0] - m0[0], m1[1] - m0[1]]
w = [inv[0][0]*dm[0] + inv[0][1]*dm[1], inv[1][0]*dm[0] + inv[1][1]*dm[1]]
proj = lambda p: w[0] * p[0] + w[1] * p[1]  # project onto w
print([round(v, 3) for v in w])
print(round(proj(m0), 2), round(proj(m1), 2))
# UPDATE: Fisher LDA direction w = Sw^-1 (m1 - m0)
\end{Verbatim}
\textbf{Output:}
\begin{Verbatim}[fontsize=\footnotesize,frame=single]
[1.194, 0.111]
2.64 8.97
\end{Verbatim}
\subsection*{Multiple linear regression}
Normal equations $X^TXw=X^Ty$. Strongly correlated features make $X^TX$ nearly singular; drop one or use ridge.
\begin{Verbatim}[fontsize=\footnotesize,frame=single]
X = [[1, 1, 2], [1, 2, 1], [1, 3, 4], [1, 4, 3], [1, 5, 5]]   # 1, x1, x2
y = [3, 4, 9, 10, 14]
n = 3
A = [[sum(r[i]*r[j] for r in X) for j in range(n)] for i in range(n)]   # X^T X
b = [sum(r[i]*t for r, t in zip(X, y)) for i in range(n)]               # X^T y
for i in range(n):                         # Gaussian elimination
    p = A[i][i]
    A[i] = [v / p for v in A[i]]; b[i] /= p
    for k in range(n):
        if k != i:
            f = A[k][i]
            A[k] = [a - f * c for a, c in zip(A[k], A[i])]; b[k] -= f * b[i]
print([round(v, 2) for v in b])            # intercept, w1, w2
# UPDATE: two features that move together make X^T X nearly singular
\end{Verbatim}
\textbf{Output:}
\begin{Verbatim}[fontsize=\footnotesize,frame=single]
[-1.0, 2.0, 1.0]
\end{Verbatim}
\section{Artificial Intelligence}
\subsection*{Iterative deepening}
Depth-limited DFS with limits $0,1,2,\dots$: memory $O(bd)$ like DFS, shallowest goal like BFS.
\begin{Verbatim}[fontsize=\footnotesize,frame=single]
g = {"A": ["B", "C"], "B": ["D", "E"], "C": ["F", "G"], "D": [], "E": [], "F": [], "G": []}
def dls(u, goal, limit, seen):             # depth-limited DFS
    seen.append(u)
    if u == goal: return True
    if limit == 0: return False
    return any(dls(v, goal, limit - 1, seen) for v in g[u])
total = 0
for depth in range(4):                     # deepen one level at a time
    seen = []
    found = dls("A", "G", depth, seen)
    total += len(seen)
    print(depth, found, len(seen))
    if found: break
print(total)
# UPDATE: memory like DFS (O(bd)), answer like BFS (shallowest)
\end{Verbatim}
\textbf{Output:}
\begin{Verbatim}[fontsize=\footnotesize,frame=single]
0 False 1
1 False 3
2 True 7
11
\end{Verbatim}
\subsection*{Conditional independence and d-separation}
Observing the middle node blocks a chain or fork but opens a collider (explaining away).
\begin{Verbatim}[fontsize=\footnotesize,frame=single]
from itertools import product
def indep(joint, given):                   # is A independent of C (given B?)
    for g in ([0, 1] if given else [None]):
        keep = {s: p for s, p in joint.items() if g is None or s[1] == g}
        z = sum(keep.values())
        pa = {a: sum(p for s, p in keep.items() if s[0] == a) / z for a in (0, 1)}
        pc = {c: sum(p for s, p in keep.items() if s[2] == c) / z for c in (0, 1)}
        for a, c in product((0, 1), repeat=2):
            pac = sum(p for s, p in keep.items() if s[0] == a and s[2] == c) / z
            if abs(pac - pa[a] * pc[c]) > 1e-9: return False
    return True
S = list(product((0, 1), repeat=3))        # state = (A, B, C)
chain = {s: .5 * (.9 if s[1] == s[0] else .1) * (.8 if s[2] == s[1] else .2) for s in S}
coll = {s: .25 * (.9 if s[1] == (s[0] | s[2]) else .1) for s in S}
print(indep(chain, False), indep(chain, True))   # chain A -> B -> C
print(indep(coll, False), indep(coll, True))     # collider A -> B <- C
# UPDATE: observing the middle blocks a chain but OPENS a collider
\end{Verbatim}
\textbf{Output:}
\begin{Verbatim}[fontsize=\footnotesize,frame=single]
False True
True False
\end{Verbatim}
\section{Database Management}
\subsection*{Relational algebra}
$\sigma$ selects rows, $\pi$ projects columns, join combines on keys, division finds ``for all'' matches; the SQL form uses GROUP BY and HAVING.
\begin{Verbatim}[fontsize=\footnotesize,frame=single]
import sqlite3
S = {("s1", "ann", "cs"), ("s2", "bob", "ee"), ("s3", "cy", "cs")}      # Student(id,name,dept)
E = {("s1", "db"), ("s1", "ml"), ("s2", "db"), ("s3", "db"), ("s3", "ml")}  # Enrolled(id,course)
sel = {r for r in S if r[2] == "cs"}                  # sigma: pick rows
proj = {(r[1],) for r in S}                           # pi: pick columns
join = {(a[0], a[1], b[1]) for a in S for b in E if a[0] == b[0]}   # natural join
print(sorted(sel), len(join))
need = {"db", "ml"}                                   # division: enrolled in ALL of need
div = {s for s, _ in E if need <= {c for t, c in E if t == s}}
print(sorted(div))
db = sqlite3.connect(":memory:")
db.execute("create table e(id, course)")
db.executemany("insert into e values (?, ?)", sorted(E))
q = "select id from e where course in ('db','ml') group by id having count(distinct course) = 2"
print(sorted(r[0] for r in db.execute(q)))            # same answer in SQL
# UPDATE: relational algebra uses SETS; SQL tables are BAGS
\end{Verbatim}
\textbf{Output:}
\begin{Verbatim}[fontsize=\footnotesize,frame=single]
[('s1', 'ann', 'cs'), ('s3', 'cy', 'cs')] 5
['s1', 's3']
['s1', 's3']
\end{Verbatim}
\subsection*{Sets versus bags}
Algebra works on sets (no duplicates); SQL tables are bags, so use DISTINCT to match $\pi$.
\begin{Verbatim}[fontsize=\footnotesize,frame=single]
import sqlite3
db = sqlite3.connect(":memory:")
db.execute("create table s(id, dept)")
db.executemany("insert into s values (?, ?)", [(1, "cs"), (2, "cs"), (3, "ee")])
print(db.execute("select dept from s").fetchall())            # SQL keeps duplicates
print(db.execute("select distinct dept from s").fetchall())   # pi removes them
print(len({"cs", "cs", "ee"}))                                # a set has no duplicates
# UPDATE: add DISTINCT to match the algebra
\end{Verbatim}
\textbf{Output:}
\begin{Verbatim}[fontsize=\footnotesize,frame=single]
[('cs',), ('cs',), ('ee',)]
[('cs',), ('ee',)]
2
\end{Verbatim}
\end{document}
